%%% solo-bloecke.tex — die Bloecke des Solos, in Autorenreihenfolge.
%%%
%%% Diese Datei wird nie unmittelbar gesetzt: werkzeuge/solo.py teilt sie
%%% bytegetreu auf und schreibt die Reihenfolge, in der sie im Heft
%%% erscheinen. Die Reihenfolge hier ist also allein die des Schreibens.
%%%
%%% Die Marken sind ASCII und ohne Leerzeichen; das Werkzeug prueft das.

\begin{soloBlock}{start}
Der Regen hat aufgehört, als du die letzte Anhöhe vor Havena erreichst. Unter dir
liegt die Stadt im fahlen Licht des späten Nachmittags, und der Geruch von Salz und
nassem Holz steigt dir entgegen. In deiner Gürteltasche steckt ein versiegelter
Brief, den du bis Sonnenuntergang übergeben sollst — an wen, das steht auf dem
Siegel, und das Siegel darfst du nicht brechen.

Zwei Wege führen hinab. Der breite endet am Stadttor, wo die Wachen jeden Ankömmling
mustern; der schmale verliert sich im Wald und kommt weiter unten wieder heraus.

Nimmst du den breiten Weg zum Tor, lies bei Abschnitt \soloWeiter{tor} weiter.
Willst du lieber ungesehen bleiben und durch den Wald gehen, so wende dich an
Abschnitt \soloWeiter{wald}.
\end{soloBlock}

\begin{soloBlock}{tor}
Das Stadttor von Havena ist so breit, dass drei Fuhrwerke nebeneinander hindurchpassen,
und heute stehen sie alle drei still. Eine Schlange von Händlern, Bauern und Pilgern
schiebt sich langsam voran. Vor dir streiten zwei Männer über den Preis eines Fasses.

Als du endlich an der Reihe bist, hebt der Wachhabende die Hand. »Dein Geschäft in
Havena?«, fragt er, und sein Blick wandert zu deiner Gürteltasche.

Willst du ihm die Wahrheit sagen, so lies bei Abschnitt \soloWeiter{wache} weiter.
Hast du ein paar Silbertaler übrig und willst sie ihm zustecken, gehe zu Abschnitt
\soloWeiter{bestechen}.
\end{soloBlock}

\begin{soloBlock}{wald}
Der schmale Pfad führt dich unter triefende Buchen. Es riecht nach Moder, und zweimal
hörst du etwas im Unterholz, das größer ist als ein Fuchs. Nach einer halben Stunde
teilt sich der Weg.

Nach links geht es an einem Felsvorsprung entlang, in dem eine Höhle klafft; nach
rechts hinunter zu einem Bach, dessen Rauschen du schon von weitem hörst.

Zur Höhle: Abschnitt \soloWeiter{hoehle}. Zum Bach: Abschnitt \soloWeiter{bach}.
\end{soloBlock}

\begin{soloBlock}{wache}
»Ein Brief«, sagst du. »Mehr weiß ich nicht.« Der Wachhabende nimmt dir die Tasche
nicht ab, aber er sieht dich lange an. »Es kommen dieser Tage viele Briefe nach
Havena«, sagt er schließlich, »und die wenigsten bringen Gutes.« Er tritt beiseite.

Du bist in der Stadt. Der Lärm schlägt über dir zusammen: Rufe, Räder, das Kreischen
von Möwen. Vor dir öffnet sich der Marktplatz, rechts führt eine enge Gasse in
Richtung Hafen.

Zum Markt: Abschnitt \soloWeiter{markt}. In die Gasse: Abschnitt \soloWeiter{gasse}.
\end{soloBlock}

\begin{soloBlock}{bestechen}
Die Münzen wechseln den Besitzer so beiläufig, dass niemand hinter dir es bemerkt.
Der Wachhabende sieht an dir vorbei und sagt laut: »Weiter, weiter, du hältst alle
auf.« Leiser fügt er hinzu: »Und geh nicht über den Markt. Da fragt heute noch
jemand anderes nach Briefen.«

Du bist drinnen. Was er gesagt hat, klingt dir noch nach, aber der Markt ist der
kürzeste Weg, und die Sonne steht schon tief.

Weiter bei Abschnitt \soloWeiter{markt}.
\end{soloBlock}

\begin{soloBlock}{hoehle}
Die Höhle ist niedriger, als sie von außen aussah, und sie riecht nach kaltem Rauch.
Jemand hat hier gelagert, und das ist noch nicht lange her. In der Asche findest du
ein Stück Siegellack — rot, wie das an deinem Brief.

Weiter hinten fällt der Boden ab. Du kannst einen schmalen Gang hinabsteigen oder
umkehren, solange es noch hell ist.

Hinabsteigen: Abschnitt \soloWeiter{keller}. Wenn du lieber zurückgehst und dabei
in etwas trittst, das sich um deinen Knöchel legt, lies Abschnitt \soloWeiter{falle}.
\end{soloBlock}

\begin{soloBlock}{bach}
Der Bach ist über die Ufer getreten, und du watest bis zu den Knien. Auf der anderen
Seite steht ein windschiefes Haus mit einem Schild, auf dem nichts mehr zu lesen ist.
Durch die Ritzen fällt Licht, und es riecht nach Zwiebeln und gebratenem Fisch.

Nach dem Regen und dem kalten Wasser wäre ein Feuer eine gute Sache.

Weiter bei Abschnitt \soloWeiter{wirtshaus}.
\end{soloBlock}

\begin{soloBlock}{markt}
Der Markt von Havena leert sich gerade. Händler schlagen ihre Stände zusammen, und
zwischen den Kisten streiten sich Möwen um Fischköpfe. Du kommst gut voran, bis dir
auffällt, dass ein Mann in einem grauen Umhang dreimal dieselbe Runde geht wie du.

Beim vierten Mal bleibt er stehen und sieht dich an.

Willst du dich in der Menge verlieren und zum Wirtshaus am Wasser durchschlagen,
lies Abschnitt \soloWeiter{wirtshaus}. Nimmst du die nächste Gasse, gehe zu
Abschnitt \soloWeiter{gasse}.
\end{soloBlock}

\begin{soloBlock}{gasse}
Die Gasse ist so eng, dass du die Wände auf beiden Seiten berühren könntest, und so
dunkel, dass du den Boden mehr ertastest als siehst. Hinter dir hallen Schritte. Sie
werden nicht schneller, und genau das ist das Beunruhigende daran.

Vor dir teilt sich der Weg: eine Treppe hinunter zum Wasser, ein Durchgang zu einem
Hof, in dem jemand eine Laterne angezündet hat.

Wenn du dich den Schritten stellst, lies Abschnitt \soloWeiter{verfolger}. Wenn du
zur Laterne gehst, wo ein Junge auf dich zu warten scheint, geht es bei Abschnitt
\soloWeiter{bote} weiter.
\end{soloBlock}

\begin{soloBlock}{wirtshaus}
Im »Anker« ist es warm, laut und voll. Du bekommst einen Platz am Ende einer Bank
und einen Krug, der nach mehr Wasser als Bier schmeckt. Neben dir erzählt ein alter
Mann jedem, der nicht schnell genug wegsieht, von einem Schiff, das gestern ohne
Mannschaft in den Hafen getrieben sei.

Die Wirtin stellt dir wortlos einen zweiten Krug hin, den du nicht bestellt hast.
»Von dem Jungen da hinten«, sagt sie. Als du dich umdrehst, ist da kein Junge mehr,
nur eine Tür, die noch nachschwingt, und dahinter eine Treppe in den Keller.

Dem Jungen nach: Abschnitt \soloWeiter{bote}. Die Kellertreppe hinunter: Abschnitt
\soloWeiter{keller}.
\end{soloBlock}

\begin{soloBlock}{keller}
Unten ist es trocken und still, und die Stille ist das Falsche daran — über dir
müsste das Wirtshaus zu hören sein. Zwischen den Fässern steht eine Kiste, die zu
sauber ist für diesen Ort.

In der Kiste liegen Briefe. Viele Briefe, alle mit rotem Siegel, alle ungeöffnet.
Deiner ist nicht der erste, der nach Havena gebracht wurde.

Nimmst du einen davon an dich, lies Abschnitt \soloWeiter{brief}. Suchst du weiter
und findest hinter der Kiste eine Truhe, so gehe zu Abschnitt \soloWeiter{schatz}.
\end{soloBlock}

\begin{soloBlock}{bote}
Der Junge kann nicht älter als zwölf sein und hat die Augen von jemandem, der lange
nicht geschlafen hat. »Du hast den Brief«, sagt er. Es ist keine Frage. »Sie haben
schon drei vor dir geschickt. Keiner ist angekommen.«

Er sieht über deine Schulter, und was er dort sieht, gefällt ihm nicht. »Nicht zum
Hafen«, sagt er schnell. »Noch nicht.«

Willst du wissen, was in dem Brief steht, lies Abschnitt \soloWeiter{brief}.
\end{soloBlock}

\begin{soloBlock}{brief}
Das Siegel bricht leichter, als es sollte — jemand hat es schon einmal geöffnet und
neu gesetzt. Der Bogen darin trägt drei Zeilen, und keine davon ergibt auf den ersten
Blick einen Sinn. Erst beim zweiten fällt dir auf, dass es keine Wörter sind, sondern
Namen. Der letzte ist deiner.

Hinter dir fällt eine Tür ins Schloss.

Wenn du dich umdrehst: Abschnitt \soloWeiter{verfolger}. Wenn du rennst und zwar
zum Hafen, wo noch Licht ist: Abschnitt \soloWeiter{hafen}.
\end{soloBlock}

\begin{soloBlock}{verfolger}
Der Mann im grauen Umhang steht drei Schritte vor dir und hat die Hände offen
erhoben, was bei jemandem mit einem solchen Gesicht nichts zu bedeuten hat. »Du
sollst mir nur den Brief geben«, sagt er. »Dann gehst du nach Hause.«

Er lügt so schlecht, dass es beinahe höflich wirkt.

Über die Dächer entkommst du vielleicht: Abschnitt \soloWeiter{dach}. Wenn du dich
an ihm vorbeidrängst und einfach läufst, lies Abschnitt \soloWeiter{flucht}.
\end{soloBlock}

\begin{soloBlock}{dach}
Die Ziegel sind nass und lose, und zweimal rutschst du so weit, dass du schon die
Gasse unter dir siehst. Aber oben ist der Wind, oben ist niemand, und von hier aus
liegt ganz Havena vor dir wie eine Karte.

Am Hafen brennen Lichter. Eines davon bewegt sich: eine Laterne an einem Mast, und
darunter macht jemand ein Schiff los.

Weiter bei Abschnitt \soloWeiter{hafen}.
\end{soloBlock}

\begin{soloBlock}{hafen}
Der Hafen von Havena riecht nach Teer, Tang und alten Netzen. Zwischen den Kais
liegen Schiffe so dicht, dass man über die Decks bis ans andere Ende laufen könnte.

An einem der hinteren Stege wird gearbeitet, obwohl es längst dunkel ist. Eine Frau
in einem schweren Mantel steht am Fallreep und zählt jeden, der an Bord geht.

Willst du unbemerkt an Bord gehen, lies Abschnitt \soloWeiter{schiff}. Sprichst du
die Frau an, so geht es bei Abschnitt \soloWeiter{kapitaen} weiter.
\end{soloBlock}

\begin{soloBlock}{schiff}
Zwischen zwei Fässern hindurch, über die Reling, hinter die Aufbauten — es ist
leichter, als es sein dürfte, und darüber solltest du dir Gedanken machen.

Unter Deck findest du einen Platz zwischen Tauen, der beinahe bequem ist. Irgendwann
löst sich das Schiff vom Steg, und das Rollen beginnt.

Weiter bei Abschnitt \soloWeiter{sturm}.
\end{soloBlock}

\begin{soloBlock}{kapitaen}
»Du bist spät«, sagt die Frau, bevor du ein Wort gesagt hast. »Und du bist nicht der,
den ich erwartet habe.« Sie hebt die Laterne und sieht dir ins Gesicht, lange.

»Aber du hast den Brief. Das genügt.« Sie tritt beiseite. »An Bord. Wir laufen aus,
bevor die Flut kippt, und ich erkläre unterwegs, wieso du besser nie in Havena
angekommen wärst.«

Gehst du an Bord: Abschnitt \soloWeiter{schiff}. Bleibst du stehen und verlangst
zuerst eine Erklärung, so lies Abschnitt \soloWeiter{sturm}.
\end{soloBlock}

\begin{soloBlock}{sturm}
Der Sturm kommt aus Südwest und kommt schnell. Das Schiff stampft, die Laterne am
Mast erlischt, und irgendwo reißt Tuch. Zweimal geht Wasser über das Deck, und beim
zweiten Mal hörst du, wie jemand über Bord geht, ohne dass du sagen könntest, wer.

Gegen Morgen legt sich der Wind so plötzlich, wie er gekommen ist. Vor dir liegt
Land, das auf keiner Karte steht.

Wenn du an Land gehst, lies Abschnitt \soloWeiter{insel}. Wenn dich eine Welle
erfasst, bevor du dich festhalten kannst, lies Abschnitt \soloWeiter{ende-tod}.
\end{soloBlock}

\begin{soloBlock}{insel}
Die Insel ist kaum größer als der Marktplatz von Havena und besteht fast nur aus
Fels. In der Mitte steht ein Turm, der dort nicht stehen kann: kein Schiff hätte
die Steine hierher gebracht, und gewachsen ist er auch nicht.

Die Tür am Fuß des Turms steht offen.

Weiter bei Abschnitt \soloWeiter{turm}.
\end{soloBlock}

\begin{soloBlock}{turm}
Im Turm ist es wärmer als draußen, und es riecht nach Staub und nach etwas Süßem,
das du nicht zuordnen kannst. Eine Treppe windet sich an der Wand nach oben. Auf
jeder dritten Stufe liegt ein Brief mit rotem Siegel.

Oben brennt Licht. Auf halber Höhe führt eine niedrige Tür in eine Kammer.

Nach oben zum Licht: Abschnitt \soloWeiter{magier}. In die Kammer: Abschnitt
\soloWeiter{kammer}.
\end{soloBlock}

\begin{soloBlock}{magier}
Der Mann am Tisch ist alt, aber nicht gebrechlich, und er schreibt, ohne aufzusehen.
»Setz dich«, sagt er. »Du bist der Vierte. Die anderen drei haben das Siegel nicht
gebrochen, deshalb sind sie noch unterwegs.« Er lächelt, und es ist kein böses
Lächeln, was die Sache schlimmer macht.

»Der Brief war nie für jemanden in Havena. Er war dafür da, dass du hierherkommst.«

Willst du wissen, was in der Kammer unter dir liegt, lies Abschnitt
\soloWeiter{kammer}. Willst du wissen, was dein Name auf dieser Liste zu suchen hat,
und bekommst die Antwort, so lies Abschnitt \soloWeiter{ende-verrat}.
\end{soloBlock}

\begin{soloBlock}{kammer}
Die Kammer ist niedrig, und du musst den Kopf einziehen. An den Wänden stehen Regale
voller Kästen, und auf jedem Kasten steht ein Name. Auf dem vordersten steht deiner.

Der Kasten ist nicht verschlossen.

Öffnest du ihn und findest darin, was du dort befürchtet hast, lies Abschnitt
\soloWeiter{schatz}. Trittst du zurück und stößt dabei gegen das Regal hinter dir,
so lies Abschnitt \soloWeiter{falle}.
\end{soloBlock}

\begin{soloBlock}{schatz}
Es ist kein Gold. In dem Kasten liegt ein Bündel Briefe, alle in deiner eigenen
Handschrift, alle mit Daten, die noch nicht gewesen sind. Der oberste beschreibt,
wie du diesen Kasten öffnest.

Du liest bis zum Ende, und am Ende steht, was du jetzt tun wirst. Und du tust es,
nicht weil es dort steht, sondern weil es das Richtige ist — und genau das, erkennst
du, ist der Unterschied, auf den es die ganze Zeit angekommen ist.

Weiter bei Abschnitt \soloWeiter{ende-gut}.
\end{soloBlock}

\begin{soloBlock}{falle}
Das Seil zieht sich zu, bevor du begreifst, dass es ein Seil ist. Du hängst kopfüber,
die Tasche fällt dir aus dem Gürtel, und der Brief liegt aufgeschlagen unter dir im
Staub, wo du ihn lesen kannst und nicht erreichen.

Schritte nähern sich, ohne Eile.

Wenn du dich losschneidest, bevor sie da sind, lies Abschnitt \soloWeiter{flucht}.
Wenn du es nicht schaffst, lies Abschnitt \soloWeiter{ende-tod}.
\end{soloBlock}

\begin{soloBlock}{flucht}
Du läufst, und du läufst gut. Zweimal biegst du falsch ab und einmal richtig, und
am Ende stehst du keuchend an einer Mauer, hinter der es nach Meer riecht.

Hinter dir ist es still geworden. Das heißt nicht, dass niemand mehr da ist.

Weiter bei Abschnitt \soloWeiter{ruecken}.
\end{soloBlock}

\begin{soloBlock}{ruecken}
An der Mauer sitzt der Junge aus der Gasse, als hätte er dort die ganze Zeit
gewartet. »Du hast ihn gelesen«, sagt er. »Das haben die anderen nicht.«

Er streckt die Hand aus, nicht fordernd, nur offen. Was du hineinlegst, entscheidet,
was aus dieser Nacht wird.

Gibst du ihm den Brief, lies Abschnitt \soloWeiter{ende-gut}. Behältst du ihn und
gehst allein weiter, so lies Abschnitt \soloWeiter{ende-verrat}.
\end{soloBlock}

\begin{soloBlock}{ende-gut}
Am Morgen liegt Havena unter einem Himmel, der zum ersten Mal seit Tagen keine Farbe
von Blei hat. Der Brief ist dort, wo er hingehört, und du bist es auch.

Was in den drei Zeilen stand, wirst du nicht vergessen. Aber du hast entschieden,
was daraus wird, und das ist mehr, als den meisten vergönnt ist, deren Namen auf
solchen Listen stehen.

\emph{Dieses Abenteuer ist zu Ende. Du hast es gut zu Ende gebracht.}
\soloEnde
\end{soloBlock}

\begin{soloBlock}{ende-tod}
Das Letzte, was du siehst, ist rot, und du hältst es für das Siegel, bis du merkst,
dass es das nicht ist.

Irgendwo in Havena wird morgen jemand einen vierten Boten losschicken.

\emph{Dieses Abenteuer ist hier zu Ende. Beginne von vorn, wenn du magst — und geh
diesmal den anderen Weg.}
\soloEnde
\end{soloBlock}

\begin{soloBlock}{ende-verrat}
Du behältst den Brief, und eine Zeit lang scheint das die kluge Entscheidung gewesen
zu sein. Du kommst weit, du kommst gut unter, und niemand fragt nach deinem Namen.

Erst viel später, in einer Stadt, die weit genug weg ist, findest du in deiner
eigenen Tasche einen Brief mit rotem Siegel, den du nie hineingelegt hast. Darauf
steht ein Name, und es ist nicht mehr deiner.

\emph{Dieses Abenteuer ist zu Ende. Es gibt Enden, die später anfangen.}
\soloEnde
\end{soloBlock}
